% Cap.~1 Fig.~1.4 --- why Markov + Monte Carlo + ML
% Layout: each column has three fixed vertical bands
%   header (y\approx2.3) | graphic (y\approx0.9) | text (y\approx-1.3)
\begin{figure}[htbp]
\centering
\resizebox{\textwidth}{!}{%
\begin{tikzpicture}[>=Stealth]

% Column frames (tall enough for three bands)
\foreach \x in {0, 5.2, 10.4} {
  \node[draw=black!40, rounded corners=3pt, thick, fill=white,
    minimum width=4.8cm, minimum height=6.4cm] at (\x,0) {};
}

% ========== MARKOV ==========
\node[font=\bfseries\small, text=teal!40!black] at (0,2.7) {Markov chain};
\node[font=\footnotesize\itshape, text=black!55] at (0,2.25)
  {What is the structure?};

% graphic band centred at y=0.95
\begin{scope}[shift={(0,0.95)}]
  \node[draw=teal!60!black, fill=teal!15, rounded corners=1pt,
    font=\scriptsize\bfseries, minimum width=0.8cm, minimum height=0.42cm]
    (m1) at (-1.25,0.25) {S$_1$};
  \node[draw=teal!60!black, fill=teal!15, rounded corners=1pt,
    font=\scriptsize\bfseries, minimum width=0.8cm, minimum height=0.42cm]
    (m2) at (0,0.25) {S$_2$};
  \node[draw=teal!60!black, fill=teal!15, rounded corners=1pt,
    font=\scriptsize\bfseries, minimum width=0.8cm, minimum height=0.42cm]
    (m3) at (1.25,0.25) {S$_3$};
  \draw[->, thick] (m1) -- (m2);
  \draw[->, thick] (m2) -- (m3);
  \draw[->, thick, red!60!black]
    (m3.south) .. controls +(0,-0.55) and +(0,-0.55) .. (m2.south);
\end{scope}

\node[font=\footnotesize, text width=4.1cm, align=center, text=black!75]
  at (0,-1.35)
  {Interpretable transition matrix.\\[2pt]
   Process backbone with rework.\\[2pt]
   Order \emph{tested}, not assumed.};

% ========== MONTE CARLO ==========
\node[font=\bfseries\small, text=teal!40!black] at (5.2,2.7) {Monte Carlo};
\node[font=\footnotesize\itshape, text=black!55] at (5.2,2.25)
  {And the uncertainty?};

\begin{scope}[shift={(5.2,0.70)}]
  \draw[thick, black!45] (-1.45,0) -- (1.45,0);
  \foreach \h/\x in {0.28/-1.0, 0.55/-0.5, 0.95/0, 0.70/0.5, 0.32/1.0} {
    \fill[teal!45!black!70] (\x-0.16,0) rectangle (\x+0.16,\h);
  }
  \node[font=\tiny, text=black!60] at (-0.85,-0.28) {P50};
  \node[font=\tiny, text=black!60] at (0.0,-0.28) {P85};
  \node[font=\tiny, text=black!60] at (0.85,-0.28) {P95};
\end{scope}

\node[font=\footnotesize, text width=4.1cm, align=center, text=black!75]
  at (5.2,-1.35)
  {Simulated trajectories.\\[2pt]
   Full distribution of dates.\\[2pt]
   Risk---not a single point.};

% ========== MACHINE LEARNING ==========
\node[font=\bfseries\small, text=teal!40!black] at (10.4,2.7) {Machine learning};
\node[font=\footnotesize\itshape, text=black!55] at (10.4,2.25)
  {What else predicts?};

\begin{scope}[shift={(10.4,0.95)}]
  \node[draw=black!40, rounded corners=1pt, fill=orange!12,
    font=\tiny, minimum width=1.45cm] (f1) at (-1.0,0.55) {churn};
  \node[draw=black!40, rounded corners=1pt, fill=orange!12,
    font=\tiny, minimum width=1.45cm] (f2) at (-1.0,0.0) {review lat.};
  \node[draw=black!40, rounded corners=1pt, fill=orange!12,
    font=\tiny, minimum width=1.45cm] (f3) at (-1.0,-0.55) {commits};
  \node[draw=teal!50!black, fill=teal!15, rounded corners=2pt,
    font=\scriptsize\bfseries, minimum width=1.15cm, minimum height=1.05cm]
    (mdl) at (1.05,0.0) {ML};
  \draw[->, thick] (f1) -- (mdl);
  \draw[->, thick] (f2) -- (mdl);
  \draw[->, thick] (f3) -- (mdl);
\end{scope}

\node[font=\footnotesize, text width=4.1cm, align=center, text=black!75]
  at (10.4,-1.55)
  {Non-linear residuals.\\[2pt]
   Tests H1: process features\\vs.\ traditional MSR.};

% ========== H3 ==========
\node[
  draw=teal!55!black, fill=teal!8, rounded corners=3pt, thick,
  text width=14.8cm, align=center, inner sep=8pt, font=\small
] at (5.2,-3.85)
  {\textbf{H3:} combination $>$ each technique alone.
   SDLC integration gap in the SLR
   (Chapter~\ref{ch:slr}; $L3=0$ / 340 studies)
   --- niche occupied by this thesis.};

\end{tikzpicture}%
}
\caption{Methodological rationale for the PATHCAST toolset.
Each technique answers a question the others do not; removing any one loses
a capability---structure (Markov), uncertainty (Monte Carlo), or residual
predictors (machine learning).}
\label{fig:cap1-three-tools}
\end{figure}
